% ===================================================================
%  TABELL Nr 6 — Huset invändigt, möblemanget, födan, belysningen.
%  Traduction 2026 d'apres texto/io, controlee sur texto/fr et
%  texto/es. Consigne : voir texto/sv/10-tabell-01.tex.
%
%  Deux notes, toutes deux marquees « (*) », et deux appels.
%  LA FEMME DE CHAMBRE PREND SON (16), contre le fac-simile ido :
%  ecart declare dans outils/renvoji.py.
% ===================================================================

%%K t06-apar-1 apar
\VUsaut{6.0mm}
\VUtitre{15.0pt}{60}{TABELL Nr 6}
\VUsaut{1.4mm}
\VUtitre{11.4pt}{0}{Huset invändigt, möblemanget, födan,}
\VUsaut{0.4mm}
\VUtitre{11.4pt}{0}{belysningen.}
\VUsaut{2.2mm}

%%K t06-apar-2 apar
Huset är färdigbyggt. Johans farbror är installerad i sin nya bostad, vars indelning vi
känner. Låt oss nu se hur han har möblerat sitt hus. Vi ska först utforska:

%%K t06-tit-1 sub
\VUpk{10.2pt}{40}{Sovrummet.}
\VUsaut{3.0mm}

%%K t06-c1-01-1 p
1. --- Vi kommer olägligt, för kammarjungfrun håller på att städa \VUgras{rummet}.

%%K t06-c1-02-1 p
2. --- På \VUgras{toalettbordet} \textsuperscript{(1)} ligger här och där olika föremål med
vilka man tvättar sig, kammar sig, kort sagt sköter sitt yttre. De är
\VUgras{tvättfatet} \textsuperscript{(2)} och \VUgras{vattenkannan}
\textsuperscript{(3)}, \VUgras{hårborsten} \textsuperscript{(4)}, \VUgras{kammen}
\textsuperscript{(5)}, \VUgras{tvålen} \textsuperscript{(6)} och
\VUgras{svampen} \textsuperscript{(7)}, slutligen en liten \VUgras{flaska}
\textsuperscript{(8)} för det flytande tandmedlet.

%%K t06-c1-03-1 p
3. --- På golvet, framför toalettbordet, här är \VUgras{hinken} \textsuperscript{(9)} för att
samla upp det smutsiga vattnet.

%%K t06-c1-04-1 p
4. --- Vi ser i \VUgras{förrummet} \textsuperscript{(10)} några
\VUgras{klädkrokar} \textsuperscript{(13)} och två \VUgras{paraplyer}
\textsuperscript{(11)} i \VUgras{paraplystället} \textsuperscript{(12)}.

%%K t06-c1-05-1 p
5. --- På andra sidan om dörren står \VUgras{spegelskåpet} \textsuperscript{(14)} som
innehåller \VUgras{linnet} \textsuperscript{(15)}.

%%K t06-c1-06-1 p
6. --- Låt oss dra nytta av att \VUgras{kammarjungfrun} \textsuperscript{(16)} bäddar
\VUgras{sängen} \textsuperscript{(17)}, för att utforska dess olika delar.
\VUgras{Sängramen} \textsuperscript{(20)} innehåller en god
\VUgras{resårmadrass} \textsuperscript{(21)} på vilken Justina genast ska lägga
en \VUgras{madrass} \textsuperscript{(22)} av ull. Sedan ska hon ta de två
\VUgras{lakanen} \textsuperscript{(23)} och \VUgras{filten}
\textsuperscript{(24)} från stolarna. Då ska hon lägga på huvudgärden den
\VUgras{tvärgående kudden} \textsuperscript{(25)} och \VUgras{huvudkudden}
\textsuperscript{(26)} som ligger med \VUgras{ejderdunstäcket} \textsuperscript{(27)} på
\VUgras{sängmattan} \textsuperscript{(32)}. Hon använder en dammvippa för att damma av
\VUgras{gardinerna} \textsuperscript{(19)} och \VUgras{sänghimlen} \textsuperscript{(18)} och
här är en del av arbetet färdigt.

%%K t06-c1-07-1 p
7. --- Då ska hon behöva ordna \VUgras{nattduksbordet} \textsuperscript{(28)} en smula, dra
upp \VUgras{väckarklockan} \textsuperscript{(29)}, göra i ordning
\VUgras{nattlampan} \textsuperscript{(30)}, ta bort \VUgras{ljuset}
\textsuperscript{(31)} för att göra rent ljusstaken.

%%K t06-tit-2 sub
\VUsaut{5.0mm}
\VUpk{10.2pt}{40}{Sykorgen och eldstadens tillbehör.}
\VUsaut{3.0mm}

%%K t06-c1-08-1 p
8. --- \VUgras{Sykorgen} \textsuperscript{(34)} bredvid \VUgras{pallen} \textsuperscript{(33)}
behöver också mycket väl ordnas. Hon ska behöva vika ihop \VUgras{broderierna}
\textsuperscript{(35)}, låsa in i \VUgras{sybordet} \textsuperscript{(38)}
\VUgras{fingerborgen}, \VUgras{tråden} \textsuperscript{(36)}, \VUgras{nålarna}
\textsuperscript{(37)} och
\VUgras{saxen} \textsuperscript{(39)} och stänga bordets lock med
\VUgras{nyckeln} \textsuperscript{(40)}.

%%K t06-c1-09-1 p
9. --- På det \VUgras{enbenta bordet} \textsuperscript{(41)} i mitten ska Justina byta vattnet
i \VUgras{karaffen} \textsuperscript{(42)}. Hon ska lägga
\VUgras{socker} i \VUgras{sockerskålen} \textsuperscript{(43)}, göra rent
\VUgras{sockertången} \textsuperscript{(44)} och ta bort \VUgras{klädborsten}
\textsuperscript{(46)}.

%%K t06-c1-10-1 p
10. --- Rummets högra sida ska kräva mindre arbete av henne, för \VUgras{schäslongen}
\textsuperscript{(47)} och dess \VUgras{kudde} \textsuperscript{(48)} står på sin rätta plats,
och eftersom det inte är kallt väder, är ingenting flyttat bland \VUgras{eldstadens}
\textsuperscript{(49)} tillbehör (redskap). \VUgras{Skoveln} \textsuperscript{(50)},
\VUgras{tången} \textsuperscript{(51)}, \VUgras{eldbockarna} \textsuperscript{(55)},
\VUgras{askgallret} \textsuperscript{(56)} är rena; \VUgras{eldskärmen}
\textsuperscript{(54)} har inte flyttats och \VUgras{vedlådan} \textsuperscript{(52)} är väl
försedd med \VUgras{ved} \textsuperscript{(53)}.

%%K t06-noto-1 noto
\VUnotes{143.2mm}{%
(*) Babo = litet barn, det E. baby, F. bébé, I. bambino.%
}

%%K t06-noto-2 noto
\VUnotes{143.2mm}{%
(*) Lambrequino = F. lambrequin, S. lambrequines, Port. lambrequins, till och med T. E.
lambrequins, alltså T. E. F. P. S.%
}

%%K t06-c1-11-1 p
11. --- Ändå ska hon behöva damma av \VUgras{pendylen} \textsuperscript{(57)},
\VUgras{kandelabrarna} \textsuperscript{(58)} och
\VUgras{småskrinen} \textsuperscript{(59)}, slutligen torka av \VUgras{spegeln}
\textsuperscript{(60)}.

%%K t06-c1-12-1 p
12. --- Vad \VUgras{vaggan} \textsuperscript{(61)} åt det lilla barnet (åt babon (*))
beträffar, står den redan färdig.

%%K t06-tit-3 sub
\VUsaut{5.0mm}
\VUpk{10.2pt}{40}{Salongen.}
\VUsaut{3.0mm}

%%K t06-c2-01-1 p
1. --- Nu här är salongen. Det är ett mycket vackert \VUgras{rum}. Eldstaden, som står i högra
hörnet är prydd med en fin \VUgras{statyett} \textsuperscript{(1)}, med två vackra
\VUgras{vaser} \textsuperscript{(3)} och draperad med en \VUgras{spiskappa}
\textsuperscript{(2)} av sammet (*).

%%K t06-c2-02-1 p
2. --- För att hindra att gnistorna från härden antänder mattan, har man ställt framför härden
ett \VUgras{gnistgaller} \textsuperscript{(4)} av koppar.

%%K t06-c2-03-1 p
3. --- Nära intill står ett elegant \VUgras{skrivbord} \textsuperscript{(5)} för damer öppet.
Vid det skriver \VUgras{husets fru} \textsuperscript{(23)} sina brev.

%%K t06-c2-04-1 p
4. --- Fönstret är försett med \VUgras{tunna gardiner} \textsuperscript{(6)} med
\VUgras{band} \textsuperscript{(8)} fästa vid \VUgras{krokar}
\textsuperscript{(7)} och med tygdraperier \VUgras{upptill draperade} \textsuperscript{(9)}.

%%K t06-c2-05-1 p
5. --- I fönsternischen innehåller en \VUgras{ytterkruka} \textsuperscript{(11)} en grön
\VUgras{växt} \textsuperscript{(10)}, en praktfull palm vars blad sträcker sig nästan ända
till \VUgras{fotogenlampan} \textsuperscript{(12)}.

%%K t06-c2-06-1 p
6. --- Lampans \VUgras{skärm} \textsuperscript{(13)} är ordnad av Maria, den förtjusande
\VUgras{unga flickan} \textsuperscript{(14)} som sitter vid pianot \textsuperscript{(15)}.
Sittande mycket rak på \VUgras{pallen} \textsuperscript{(16)}, övar hon ett musikstycke. Hon
är mycket skicklig och omsorgsfull, som hennes \VUgras{notskåp} \textsuperscript{(17)} bevisar
som är fullkomligt ordnat.

%%K t06-tit-4 sub
\VUsaut{5.0mm}
\VUpk{10.2pt}{40}{Odågan.}
\VUsaut{3.0mm}

%%K t06-c2-07-1 p
7. --- Hennes bror, Paul, söker en \VUgras{bok} \textsuperscript{(19)} i
\VUgras{bokhyllan} \textsuperscript{(18)}. Han är en \VUgras{ung pojke}
\textsuperscript{(20)}, rörlig och odräglig. Medan han hänger sig fast vid bokhyllan för att
nå boken som står för högt, riskerar han att fälla ner
\VUgras{bysten} \textsuperscript{(21)} som står ovanpå.

%%K t06-c2-08-1 p
8. --- Han begagnar sig av att hans mor är upptagen med att samtala med en väninna, en
\VUgras{dam} \textsuperscript{(22)} som kom för att tillbringa några dagar hos henne, för att
göra den dumheten. Modern sitter på en
\VUgras{soffa} \textsuperscript{(24)} bredvid \VUgras{länstolen}
\textsuperscript{(25)}, och hennes väninna sitter på \VUgras{stolen} \textsuperscript{(27)} av
vilken man bara ser \VUgras{ryggstödet} \textsuperscript{(26)}.

%%K t06-c2-09-1 p
9. --- Den stora \VUgras{ramen} \textsuperscript{(30)} som hänger på väggen innehåller
\VUgras{porträttet} \textsuperscript{(29)} av en släkting. I
\VUgras{albumet} \textsuperscript{(32)} som ligger på bordet är de övriga
familjemedlemmarnas \VUgras{fotografier} \textsuperscript{(33)} samlade. Barnen har just
bläddrat i det, för \VUgras{stolarna} \textsuperscript{(31)} står ännu grupperade kring
bordet.

%%K t06-c2-10-1 p
10. --- Den \VUgras{kinesiska vasen} \textsuperscript{(34)} av porslin, som står bredvid
\VUgras{papperskniven} \textsuperscript{(35)} skänktes åt Maria till hennes namnsdag, och
\VUgras{skärmen} \textsuperscript{(36)} som pryder rummets hörn fördes hem från Kina av hennes
farbror.

%%K t06-c2-11-1 p
11. --- Upplivad av de praktfulla \VUgras{målningarna} \textsuperscript{(37)} som hänger på
\VUgras{väggen} \textsuperscript{(42)}, upplyst av
\VUgras{takkronan} \textsuperscript{(38)}, vars \VUgras{elektriska lampor}
\textsuperscript{(39)} lyser glatt om kvällen, ser den salongen mycket behaglig ut. Den mjuka
\VUgras{mattan} \textsuperscript{(40)} som ligger utbredd på parketten, och den så lättanvända
\VUgras{elektriska ringklockan} \textsuperscript{(41)}, fullbordar dess bekvämlighet.

%%K t06-tit-5 sub
\VUsaut{3.6mm}
\VUpk{10.2pt}{40}{Matsalen.}
\VUsaut{3.0mm}

%%K t06-c3-01-1 p
1. --- Middagsklockan har ringt. Hela familjen, utom Maria, är samlad kring bordet. Matsalen
värms behagligt av den utmärkta \VUgras{kaminen} \textsuperscript{(1)} av fajans, med ett tunt
\VUgras{rör} \textsuperscript{(2)}.

%%K t06-c3-02-1 p
2. --- Det rummet innehåller inte många möbler. Längs väggen hänger, ovanför
\VUgras{bänken} \textsuperscript{(4)}, en prydlig \VUgras{fontän}
\textsuperscript{(3)}. Alldeles intill står \VUgras{skänken} \textsuperscript{(5)} på vilken
man ställer den kalla maten, efterrättsfaten, och de olika bordsredskap som man inte behöver
genast.

%%K t06-c3-03-1 p
3. --- Så ser vi på den övre hyllan en \VUgras{stekt kyckling} \textsuperscript{(7)}, en
\VUgras{fisk} \textsuperscript{(8)} och en
\VUgras{skål} \textsuperscript{(10)} med \VUgras{frukt} \textsuperscript{(9)}.
Under här är \VUgras{osten} \textsuperscript{(11)}, \VUgras{brödet} \textsuperscript{(12)},
\VUgras{bordsstället} \textsuperscript{(13)} som innehåller allt vad man behöver för att
krydda salladen: oljan, ättikan,
\VUgras{saltet} \textsuperscript{(14)} och \VUgras{pepparen}
\textsuperscript{(15)}. Slutligen står på den nedre hyllan en \VUgras{tårta}
\textsuperscript{(17)} bredvid \VUgras{senapsburken} \textsuperscript{(16)}.

%%K t06-c3-04-1 p
4. --- Matsalens ur, som man kallar \VUgras{väggur} \textsuperscript{(20)}, har en mycket
säregen form. Det är infattat i en träram som dessutom innehåller en
\VUgras{barometer} \textsuperscript{(19)} och en \VUgras{termometer}
\textsuperscript{(18)}.

%%K t06-tit-6 sub
\VUsaut{4.6mm}
\VUpk{10.2pt}{40}{Middagens servering.}
\VUsaut{3.0mm}

%%K t06-c3-05-1 p
5. --- På alla rummets väggar är \VUgras{träpaneler} \textsuperscript{(21)} av vaxad ek
anbragta. En \VUgras{portiär} \textsuperscript{(22)} av tyg stänger tillräckligt dörren som
ger tillträde till verandan. Dit ska familjen gå för att dricka kaffet efter middagen. Redan
står \VUgras{kopparna} \textsuperscript{(24)} och \VUgras{faten} \textsuperscript{(25)}
färdiga på
\VUgras{serveringsbordet} \textsuperscript{(23)}.

%%K t06-c3-06-1 p
6. --- Ovanför bordet är en \VUgras{hänglampa} \textsuperscript{(26)} fäst vid taket vars
mycket klara lampa om kvällen lyser upp hela matsalen.

%%K t06-c3-07-1 p
7. --- Låt oss nu utforska själva \VUgras{matbordet} \textsuperscript{(27)}. När familjen kom
in, var kuvertet färdigt. På \VUgras{bordduken} \textsuperscript{(28)} var på varje bordsgästs
plats en \VUgras{tallrik} \textsuperscript{(33)}, en \VUgras{servett} \textsuperscript{(29)},
en
\VUgras{sked} \textsuperscript{(34)}, en \VUgras{gaffel}
\textsuperscript{(35)}, en \VUgras{kniv} \textsuperscript{(36)} och ett
\VUgras{glas} \textsuperscript{(32)} lagda. Det fanns också \VUgras{flaskor}
\textsuperscript{(30)} för vinet och en \VUgras{karaff} \textsuperscript{(31)} för vattnet.

%%K t06-c3-08-1 p
8. --- \VUgras{Betjänten} \textsuperscript{(39)} bar först in
\VUgras{soppskålen} \textsuperscript{(37)}, i vilken det ännu ryker litet soppa.
Nu serverar han den första \VUgras{rätten} \textsuperscript{(38)}, en kötträtt. Sedan ska han
bjuda dem som vi redan har nämnt, slutligen, efter
\VUgras{efterrättsfatet}, ska han begagna sig av att hans herrskap har gått in
i verandan, för att ta bort bordsredskapen och ordna matsalen igen.

%%K t06-c3-08-2 p
\textit{Anm.} --- \textit{De gängse orden som rör födan anges här mycket kortfattat. Listan
ska fullständigas på de två tabellerna « Hotellet » nr 13) och « Marknaden » (nr 15).}

%%K t06-tit-7 sub
\VUsaut{4.6mm}
\VUpk{10.2pt}{40}{Köket.}
\VUsaut{3.0mm}

%%K t06-c4-01-1 p
1. --- I köket, som vi börjar betrakta genom den halvöppna dörren, ser vi en pittoresk
oordning. Framför \VUgras{härden} \textsuperscript{(1)} i vilken en
\VUgras{eld} \textsuperscript{(3)} sprakar, står ett \VUgras{stekspett}
\textsuperscript{(5)}. En vacker \VUgras{stek} \textsuperscript{(7)} är
\VUgras{uppträdd på spettet} \textsuperscript{(6)} och blir färdig.

%%K t06-c4-02-1 p
2. --- \VUgras{Blåsbälgen} \textsuperscript{(8)} och \VUgras{kolskoveln}
\textsuperscript{(9)}, som man just har använt, har blivit liggande på golvet liksom
\VUgras{veden} \textsuperscript{(10)}.

%%K t06-c4-03-1 p
3. --- Spiselkransen är ordnad; \VUgras{lyktan} \textsuperscript{(11)},
\VUgras{tändsticksasken} \textsuperscript{(13)}, i vilken \VUgras{tändstickorna}
\textsuperscript{(12)} ligger, \VUgras{ljusstaken} \textsuperscript{(14)} och
\VUgras{handljusstaken} \textsuperscript{(15)} med ett \VUgras{ljus}
\textsuperscript{(16)} står på sin plats. Men den stora \VUgras{kolhinken}
\textsuperscript{(18)}, full av \VUgras{stenkol} \textsuperscript{(17)} som
\VUgras{kokerskan} \textsuperscript{(23)} just har fyllt i källaren, har blivit
stående mitt i köket.

%%K t06-c4-04-1 p
4. --- Sannerligen, den stackars kvinnan måste skynda sig för att frukosten ska vara färdig i
rätt tid. Nu står hon vid \VUgras{spisen} \textsuperscript{(19)} och rör i en \VUgras{sås}
\textsuperscript{(24)} som inte får brännas vid för mycket.

%%K t06-c4-05-1 p
5. --- I \VUgras{ugnen} \textsuperscript{(20)} gräddas en tårta, som hon har gjort i ordning
till barnens mellanmål. Ovanför spisen hänger \VUgras{sleven} \textsuperscript{(21)} och
\VUgras{skumsleven} \textsuperscript{(22)}.

%%K t06-tit-8 sub
\VUsaut{4.0mm}
\VUpk{10.2pt}{40}{Husgeråden.}
\VUsaut{3.0mm}

%%K t06-c4-06-1 p
6. --- \VUgras{Kokkärlen} \textsuperscript{(25)}, som hänger längst in i rummet, är mycket
rena. De vackra \VUgras{kastrullerna} av koppar \textsuperscript{(26)}, \VUgras{mjölkkannan}
\textsuperscript{(27)}, det lilla
\VUgras{durkslaget} \textsuperscript{(28)} och \VUgras{rivjärnet}
\textsuperscript{(29)} har blivit omsorgsfullt rengjorda.

%%K t06-c4-07-1 p
7. --- \VUgras{Saltasken} \textsuperscript{(30)} står på skänken bredvid
\VUgras{tekannan} \textsuperscript{(31)}, och den \VUgras{stora oljeflaskan}
\textsuperscript{(32)}. \VUgras{Lådan} \textsuperscript{(33)} innehåller troligen matbestick
och kökets knivar.

%%K t06-c4-08-1 p
8. --- \VUgras{Kvasten} \textsuperscript{(34)} och \VUgras{dammvippan} \textsuperscript{(35)}
står i ett hörn. Kokerskan har just använt
\VUgras{huggkubben} \textsuperscript{(36)}, hon har lämnat den vid fönstret.

%%K t06-c4-09-1 p
9. --- En \VUgras{handduk} \textsuperscript{(37)} hänger på muren, bredvid
\VUgras{tratten} \textsuperscript{(38)}. I den delen av köket är
\VUgras{halstret} \textsuperscript{(39)}, \VUgras{hackkniven}
\textsuperscript{(40)} och \VUgras{hackbrädet} \textsuperscript{(41)} undanstuvade. Sedan
står, ovanför hackkniven, på en \VUgras{hylla} \textsuperscript{(42)}, \VUgras{krukor}
\textsuperscript{(43)},
\VUgras{kokgrytan} \textsuperscript{(44)} med sitt \VUgras{lock}
\textsuperscript{(45)} och \VUgras{kittelen} \textsuperscript{(46)}.
\VUgras{Stekpannan} \textsuperscript{(47)} hänger nära intill.

%%K t06-c4-10-1 p
10. --- Bara \VUgras{kranen} av koppar \textsuperscript{(48)} kastar en glans i det hörnet.
\VUgras{Vasken} \textsuperscript{(49)} på vilken den tjocka
\VUgras{stora kannan} \textsuperscript{(50)} står, är troligen mycket bekväm med
sitt lilla skåp, i vilket man förvarar \VUgras{ättiksfatet} \textsuperscript{(51)} som är
försett med sin \VUgras{tapp} \textsuperscript{(52)}, \VUgras{avfallslådan}
\textsuperscript{(53)} och de
\VUgras{smutsiga tallrikarna} \textsuperscript{(54)}.

%%K t06-tit-9 sub
\VUsaut{4.0mm}
\VUpk{10.2pt}{40}{Andra föremål i köket.}
\VUsaut{3.0mm}

%%K t06-c4-11-1 p
11. --- \VUgras{Baljan} \textsuperscript{(55)} i vilken man tvättar
\VUgras{grönsakerna} \textsuperscript{(58)} och \VUgras{grytan}
\textsuperscript{(56)} av gjutjärn som står på \VUgras{stengolvet} \textsuperscript{(57)} har
troligen också sin plats i det skåpet.

%%K t06-c4-12-1 p
12. --- Kokerskan saknar visst inte spisar, för här är dessutom, på bordets hörn, ett
\VUgras{gaskök} \textsuperscript{(59)} på vilket vatten värms i en liten kastrull.
\VUgras{Ångan} \textsuperscript{(60)} som slipper ut ur det visar oss, att det troligen kokar.
Antagligen ska det vattnet användas till kaffet, för här är på bordets andra ände,
\VUgras{kaffepannan} \textsuperscript{(64)} och
\VUgras{kaffekvarnen} \textsuperscript{(63)}.

%%K t06-c4-13-1 p
13. --- Köket är förbundet med \VUgras{gasbrännaren} \textsuperscript{(62)} genom en lång
\VUgras{gummislang} \textsuperscript{(61)}.

%%K t06-c4-14-1 p
14. --- \VUgras{Dagligvarahjälpen} \textsuperscript{(66)} som kommer för att hjälpa kokerskan,
gör i ordning grönsakerna till middagen. När de blir rensade och tvättade, ska hon lägga dem i
\VUgras{bunken} \textsuperscript{(65)} i väntan på timmen att koka dem.

%%K t06-tit-10 sub
\VUsaut{4.0mm}
{\centering\fontsize{10.2pt}{10.2pt}\selectfont\textls[40]{\scshape Badrummet.} \textit{(Badkammaren.)}\par}
\VUsaut{3.0mm}

%%K t06-c5-01-1 p
1. --- Oss återstår bara en indiskretion att begå för att lära känna våningens förnämsta rum:
att se badrummets inredning. Det ska inte kräva lång tid, för det är inte stort, och vi ska
snabbt veta att \VUgras{badkaret} \textsuperscript{(1)} har en god \VUgras{badvärmare}
\textsuperscript{(2)}, att
\VUgras{tvålkoppen} \textsuperscript{(3)} är inom räckhåll för badarens hand och
att \VUgras{handduksstället} \textsuperscript{(5)} säkert ska låta
\VUgras{handdukarna} \textsuperscript{(4)} torka lätt.

%%K t06-c5-02-1 p
2. --- Det rummet har sina väggar täckta med \VUgras{kakel} \textsuperscript{(6)}, som gjorde
det möjligt att sätta upp en
\VUgras{duschapparat} \textsuperscript{(7)} utan att på något sätt frukta att
vattnet, som stänker tillbaka vid dess bruk, skulle skada murarna. Ett litet
\VUgras{fat} \textsuperscript{(8)} av zink som tjänar till fotbaden, gör
möblemanget fullständigt.
\VUsaut{6.0mm}
\VUfilet{22mm}
